\specialchapter{How This Book Was Made}{Colophon}{34}

\addsec{Roles}

\emph{The Hydrogen Atoms of Deep Learning} was written by Claude Opus 5.5, an AI model made by Anthropic, in October 2026. Xiaoyu Li conceived the series \emph{The Missing Textbooks} and proposed this volume: its subject, its gap (solvable models of deep learning scattered across papers in incompatible notation, and rarely accompanied by an account of where they fail), its format (a phenomenon, a minimal solvable model, a theorem with a complete proof, and a section on what the model gets wrong), a first list of models, and its position relative to existing books. Xiaoyu Li also set the direction of the series (language, design and credit) and is responsible for the book's publication.

\addsec{Process}

A written plan fixed the thesis, the sixteen models, the global notation, and a specification for each chapter, including the editor's plan for each main theorem. A style guide fixed the anatomy of a chapter, the voice, a proof style that follows Halmos, the three-way distinction between proofs, sketches and heuristic derivations, and the rule that every chapter compares an exact prediction with a simulation. Separate instances of the same model then drafted the chapters, at most five at a time, each re-deriving its results, checking them numerically, verifying its references with web search, and writing the scripts behind its figures. Five chapters were begun by one set of instances and completed by another; two of these received a separate pass to finish and check the draft before refereeing. A coordinating instance wrote the prologue, preface, notation, epilogue and this colophon, designed the volume's additions to the series typography, and integrated the whole.

Each chapter, the prologue and the appendix then received one independent referee pass from a separate instance, which re-derived the proofs, re-ran the figure scripts, checked constants and exercises numerically, and checked references and historical claims by web search. The referees corrected what they found directly. They found few errors in the main theorems, and several gaps in proofs, which they closed or marked; most of their corrections were to numbers quoted from simulations, to the statements and answers of exercises, and to claims about what other papers prove, many of which were softened or removed because they could not be confirmed. The epilogue was written last and was not refereed.

\addsec{Limits}

No human has yet checked every proof. Several results are presented as heuristic derivations, marked with $\Diamond$, because no proof is known to the author or because a proof would not fit; they are checked against simulations but not proved. Web search was available while the book was written, but the full texts of most papers could not be opened, so references were checked against search results, abstracts and secondary accounts. Some historical details, and some translations of the notation of source papers, rest on these or on the author's memory, and the chapters say so where it matters. One referee pass per chapter finds many errors but not all of them. Results from 2024 to 2026 are cited from preprints where no refereed version could be found, and are marked as such where a claim rests on them. Corrections are welcome and will be recorded in the repository of \emph{The Missing Textbooks}.

\medskip
{\small\noindent Typeset with pdf\LaTeX{} in \texttt{missingbook}, the series class built on \textsc{koma}-Script, with the volume style \texttt{hydrogen}; Palatino (newpx) for text and mathematics, biblatex and biber for references, Ti\emph{k}Z and pgfplots for figures, whose data were computed with Python, NumPy, SciPy and PyTorch. The lines on the title page are the Balmer series of hydrogen, drawn at the wavelengths given by Rydberg's formula; the wave at the foot of each opening page rises from chapter to chapter.\par}
